% ECONOMICS_PROBLEM: {"id": "C78-554", "title": "Realizability of every admissible input-block shape", "jel_code": "C78", "source_id": "OP-0207"}
% !TeX program = xelatex
\documentclass[11pt,a4paper]{article}
\usepackage[margin=23mm]{geometry}
\usepackage{fontspec}
\setmainfont{Latin Modern Roman}
\usepackage{amsmath,amssymb,mathtools}
\usepackage{xcolor,enumitem,booktabs,longtable,array,xurl,hyperref}
\definecolor{muted}{HTML}{53636C}
\hypersetup{colorlinks=true,urlcolor=blue,linkcolor=blue}
\setlength{\parindent}{0pt}
\setlength{\parskip}{5pt}
\newcommand{\R}{\mathbb R}
\newcommand{\E}{\mathbb E}
\newcommand{\N}{\mathbb N}
\newcommand{\ind}{\mathbf 1}
\newcommand{\Var}{\operatorname{Var}}
\newcommand{\Cov}{\operatorname{Cov}}
\newcommand{\supp}{\operatorname{supp}}
\DeclareMathOperator*{\argmin}{arg\,min}
\DeclareMathOperator*{\argmax}{arg\,max}
\newcommand{\entrymeta}[3]{{\sffamily\small\color{muted}\textbf{\texttt{#1}}\quad|\quad #2\par #3\par}}
\newcommand{\Problemref}[1]{\texttt{#1}}
\newcommand{\JELref}[2]{\texttt{#2} (JEL \texttt{#1})}
\begin{document}
\section*{Realizability of every admissible input-block shape}
\textbf{Problem ID:} \texttt{C78-554}\quad \textbf{JEL:} \texttt{C78}\par
% BEGIN ECONOMICS STATEMENT
\label{problem:OP-0207}
% GAME_THEORY_PROBLEM: {"local_id": "GT-077", "original_number": 77, "kind": "P", "entry_kind": "statement_only", "published": false, "review_required": true, "category": "game_theory"}
\label{game:GT-077}
{\sffamily\small\color{muted}
\noindent\textbf{ID: GT-077.}\quad\textbf{Category:} Stable matching decomposition.
\quad\textbf{Kind: P.}\quad\textbf{Status:} Conjecture in the cited source.
\par}
\subsection*{Definitions and mathematical statement}
Let \(I\) have proposer and receiver sets \(P,R\) of size \(n\geq 1\),
with strict complete rankings. An \emph{input block} is a pair
\((A,B)\), where \(\varnothing\ne A\subseteq P\),
\(B\subseteq R\), and \(|A|=|B|=k\), such that
\[
 \begin{aligned}
 \{\text{the first }k\text{ receivers in }p\text{'s ranking}\}
     &=B &&(p\in A),\\
 \{\text{the first }k\text{ proposers in }r\text{'s ranking}\}
     &=A &&(r\in B).
 \end{aligned}
\]
Let \(\mathcal B(I)\) be the set of all input blocks, including
\((P,R)\). Order it by
\((A,B)\leq(A',B')\) iff \(A\subseteq A'\) and \(B\subseteq B'\).
Blocks are laminar: intersecting blocks are comparable.
The containment Hasse diagram is therefore a rooted tree after the
whole-instance root is included. Label a block by its size \(k\).

An \emph{admissible shape of order \(n\)} is a finite rooted unordered
tree \(T\), with positive integer labels \(s(v)\), satisfying
\[
 s(\mathrm{root})=n,\qquad
 s(w)<s(v)\text{ for each child }w\text{ of }v,\qquad
 \sum_{w:\,w\text{ child of }v}s(w)\leq s(v).
\]

\paragraph{Conjecture.}
For every \(n\geq1\) and every admissible shape \(T\) of order \(n\),
there is a strict complete order-\(n\) instance \(I\) such that the
containment Hasse diagram of \(\mathcal B(I)\) is isomorphic to \(T\),
with the isomorphism preserving the root and all size labels.

\subsection*{Short English statement}
Can every shape satisfying the block-size constraints occur without
creating additional blocks?

\subsection*{Variants and scope boundaries}
All input blocks must be accounted for; exhibiting only selected blocks
does not prove realizability. A proper block need not have a complementary
block, so a block is not the same as a factor of an input-independent
partition. With the whole-instance block included, the source's
block forest is a rooted tree.

\subsection*{Source relationship and status}
This is [S1, Section 7, problem D]. The top-rank definition is the
source's input-block notion, and the admissibility constraints are those
in that problem. The source reports small-order realizations, while
accidental extra blocks remain the obstruction to a general construction.

\subsection*{References}
\begingroup\footnotesize\raggedright
\noindent[S1] Yoshiteru Ishida,
\emph{Prime Factorisation of Stable-Matching Instances: Uniqueness,
Simultaneous Products, and an Exact Census}, 2026,
arXiv:2610.05617v1.
\url{https://arxiv.org/html/2610.05617v1}.
Locators: Section 3, Theorem 9.4 (block condition) and
Proposition 9.5 (laminarity); Section 7, problem D.

% Audit: 66 corrects approval-count/general-score conflation and a duplicate.
% Audit: 69 gives a padding correction; 70 is an editorial structural agenda.
% Audit: 73 replaces the obsolete upper bound 6 by the proved upper bound 5.
% Audit: 74 distinguishes exhaustive minima from targeted-search upper bounds.
% Scope: "open" refers to the explicitly cited retrieved version, not a
% claim that all original catalogue entries remain open as of every date.


% Fragment GT-078--GT-087. Source audit: 8 October 2026.
% Bibliographic convention: Cardaliaguet--Delarue is the ICM2022 chapter,
% copyright 2022, published in the EMS proceedings volume (2023).
\par\endgroup

\subsection*{Common conventions: Source mathematical conventions}

\textbf{Finite games.} For a finite set $A$, $\Delta(A)$ is its probability
simplex. Players choose independent mixed strategies in Nash equilibrium;
correlation is allowed only when explicitly part of the solution concept.
In a game with payoff functions $u_i$ and strategy profile $x$, an additive
$\varepsilon$ Nash equilibrium satisfies
\[
 u_i(x)\geq u_i(x_i',x_{-i})-\varepsilon
 \quad\text{for every player $i$ and every admissible deviation $x_i'$}.
\]
For numerical approximation, payoffs are normalized as locally specified.
An $\varepsilon$ payoff residual is distinct from distance at most
$\varepsilon$ to the set of exact equilibria. Point convergence, convergence
of empirical play, and convergence in distance to an equilibrium set are
also distinct.

\textbf{Computational representations.} Unless stated otherwise, a complexity
question takes rational data with binary encoding; polynomial time is measured
in total bit length. A fixed-accuracy polynomial algorithm is not necessarily
polynomial in $1/\varepsilon$, and neither is necessarily polynomial in
$\log(1/\varepsilon)$. Succinct games and value, demand, or sampling oracles
must declare their interfaces. Exact real solutions require an explicit
representation; an unspecified list of arbitrary real numbers is not a
finite computational output. A claimed algorithmic barrier is meaningful
only for its specified model and reductions.

\textbf{Dynamic games.} Histories, information, observation, and allowable
randomization are part of the model. A stationary strategy depends only on
the current state; a positional strategy in a deterministic graph game
chooses one action at each controlled vertex. Nash equilibrium at an initial
state and equilibrium after every history are different requirements.
An expectation of a pathwise $\liminf$ payoff, a $\liminf$ of expected averages,
a limiting expected average, and a uniform finite-horizon equilibrium are
different criteria. None is silently substituted for another.

\textbf{Allocation and mechanisms.} A complete allocation assigns every item
exactly once unless otherwise stated. Zero-valued goods, negative values,
capacity constraints, feasibility, lotteries, and copies of goods affect
fairness definitions and are specified locally. Dominant-strategy incentive
compatibility, Bayesian incentive compatibility, universal truthfulness,
and truthfulness in expectation impose different inequalities. Whenever
an objective ratio has zero denominator, its convention is stated or those
instances are excluded.

\textbf{Infinite-dimensional models.} Expectations, integrals, stochastic
equations, and optimization objectives require their local regularity,
integrability, filtrations, and admissible domains. A backward--forward
mean-field system is a boundary-value problem; it is not an initial-value
semiflow without an additional construction. Long-time, large-population,
and vanishing-noise limits require an order of limits and a convergence
topology. If a minimum need not be attained, the objective is its infimum
or supremum and the approximation requirement is stated separately.
% END ECONOMICS STATEMENT
\end{document}
